% Narrativa generada automaticamente desde comparative_stats (M2): NO editar a mano.
% Cada cifra esta interpolada desde compute_comparative; regenerar actualiza el texto.
En pregrado (5.821 participantes frente a 134.601 no participantes, tasa base 4,15\%), el análisis prioriza el tamaño del efecto sobre la mera significancia estadística, inevitable con muestras tan grandes. Las diferencias de mayor magnitud entre grupos aparecen en número de opciones inscritas (d = 1,07, efecto grande), créditos aprobados (total) (d = 0,76, efecto mediano) y percentil de pga (programa) (d = 0,75, efecto mediano); todas favorecen a quienes participan. La asociación con el estrato social es de magnitud pequeño (V de Cramér = 0,15) (estrato 1 con 0,42\% y estrato 6 con 9,48\%): persiste un gradiente socioeconómico que debe leerse como brecha de equidad, no como criterio de selección. La asociación con el género es trivial (V de Cramér = 0,02). Todas las relaciones son asociaciones observadas, no relaciones causales; el modelo logístico multivariado se reporta como aproximación (sin p-valores analíticos) y complementa, sin sustituir, la lectura univariada.
